\subsection{All profinite coefficient rings}
\label{C:profinite-subsection}

The finite free integral coalgebras of
Lemma~\ref{C:integral-orthogonal-coalgebras} also remove
the Noetherian and local restrictions. The correct
uniform bound here is the size of an ambient idempotent
matrix, since the modules need not be free.

\begin{lemma}[Proved: bounded polynomial records for projective equivariant forms]
\label{C:profinite-projective-data}
Fix $B,D\geq0$ and a commutative ring $A$ with $2$ invertible.
A nonsingular symmetric rational $O_{m,A}$-representation
on an idempotent summand of $A^B$, with coefficient bound
$D$, has a finite polynomial record
\[
               (e,\rho,q,q^\sharp),                     \tag{C.3.PF1}
\]
where $e,q,q^\sharp$ are $B$-square matrices over $A$
and $\rho$ is a $B$-square matrix with entries in $C_D(A)$.
Its exact defining equations are
\[
\begin{gathered}
 e^2=e,\qquad \rho=e\rho e,\qquad
 \varepsilon(\rho)=e,\qquad
 \Delta(\rho_{ij})=\sum_k\rho_{ik}\otimes\rho_{kj},\\
 q=e^tqe,\qquad q^t=q,\qquad
 q^\sharp=e q^\sharp e^t,\qquad
 q^\sharp q=e,\qquad qq^\sharp=e^t,\qquad
                         \rho^t q\rho=q .
\end{gathered}                                           \tag{C.3.PF2}
\]
The final identity is in $C_{2D}(A)$, with its right
side multiplied by the unit function. All tensor
identities in this display use the fixed finite free
bases from that lemma.
These records and the finite polynomial equations for
isometries commute with every base change.
Every actual finite projective equivariant symmetric
form has such a record for some $B,D$.
If $A$ is finite, the records of fixed $B,D$ form a
finite set with a terminating exhaustive enumeration.
\end{lemma}
\begin{proof}
For $P=eA^B$ the dual is identified with $e^tA^B$
under the ordinary evaluation pairing.
A coaction on $P$ extends by zero on the complement
of $e$ to a matrix $\rho=e\rho e$ on $A^B$.
Its counit is $e$ and its coassociativity is the
displayed matrix equation. Conversely these
equations restrict to a coaction on $P$.
Every entry of a rational coaction is a finite sum
of coordinate functions, so exhaustiveness of
the $C_D$ gives some common coefficient bound.

A map $q:P\to P^*$ is represented by a matrix
$q=e^tqe$. Symmetry is exactly $q^t=q$.
It is an isomorphism precisely when it has an
inverse represented by $q^\sharp=e q^\sharp e^t$
with the two stated inverse identities.
The last identity is precisely invariance of the
associated bilinear pairing. It is a finite
coefficient identity because multiplication maps
$C_D C_D$ to $C_{2D}$.
Thus (C.3.PF2) is both necessary and sufficient.
It also covers the zero module.

For clarity, an isometry from a record on $eA^B$
to one on $fA^{B'}$, with coaction $\sigma$ and
Gram matrix $h$, is encoded by two matrices
$g\in M_{B'\times B}(A)$ and
$u\in M_{B\times B'}(A)$ satisfying
\[
\begin{gathered}
 g=fge,\qquad u=euf,\qquad ug=e,\qquad gu=f,\\
                 g\rho=\sigma g,\qquad g^t h g=q .
\end{gathered}                                           \tag{C.3.PF3}
\]
The first line makes $g,u$ inverse maps between
the indicated projectives; the second gives
equivariance and preservation of the form.
These are again finite polynomial coefficient
equations. All equations commute with base change.
Every finite projective has an idempotent
presentation, so the records exhaust the category.
Over a finite ring all matrix and coefficient
sets being enumerated are finite, proving the
last assertion.
\end{proof}

\begin{theorem}[Proved: complete bounded finite-coordinate classification for every profinite ring]
\label{C:profinite-WG}
Let
\[
        R=\varprojlim_{\alpha\in\Lambda} R_\alpha
\]
be a directed inverse limit of finite commutative
rings with $2$ invertible and surjective transition
maps. Thus $R$ is an arbitrary profinite commutative
ring with $2$ invertible, in a chosen finite quotient
presentation. Put
$G_\alpha=WG_{\rm sp}(BO_{m,R_\alpha})$, in the
explicit canonical coordinates of
Theorem~\ref{C:Om-finite-base-algorithm}.

There is an injective homomorphism preserving all
exterior operations,
\[
 WG_{\rm sp}(BO_{m,R})\longrightarrow
                       \varprojlim_\alpha G_\alpha.     \tag{C.3.PF4}
\]
Its image is exactly the set of compatible canonical
coordinate families $(\xi_\alpha)$ satisfying
\[
\begin{split}
 \exists B,D\ \forall\alpha:\quad
 \xi_\alpha\text{ is the difference of the classes}\\
 \text{of two records (C.3.PF1)--(C.3.PF2)
 over }R_\alpha\text{ of the same bounds }B,D .
\end{split}                                             \tag{C.3.PF5}
\]
For fixed $B,D,\alpha$, this is a finite decidable
test. All transition maps, products and specified
exterior coefficients in these coordinates are
computed by the finite-base algorithms.
Hence (C.3.PF4)--(C.3.PF5) give the entire split ring
as a literal bounded subset of compatible explicitly
solved finite-coordinate groups, without an
unclassified isomorphism quotient.
No ordinary finite decision procedure for equality
of arbitrary infinite families is asserted.

This includes genuinely non-Noetherian local
profinite rings, for example
\[
 R=\mathbf F_p\ltimes\prod_{j\geq1}\mathbf F_p\epsilon_j,
 \qquad
 \left(\prod_{j\geq1}\mathbf F_p\epsilon_j\right)^2=0,
 \qquad p>2,                                           \tag{C.3.PF6}
\]
with the topology of truncation in the $\epsilon_j$.
\end{theorem}
\begin{proof}
We use the following elementary compactness fact.
An inverse system of nonempty finite sets indexed
by a directed set has a nonempty inverse limit,
even when its individual transition maps are not
surjective. Indeed, put the discrete topology on
the finite sets. The compatibility conditions
are closed in their compact product.
Any finite collection of those conditions mentions
only finitely many indices. Choose one upper bound,
choose an element at that level and push it to
all mentioned indices. Thus every finite collection
is consistent, and compactness proves the assertion.
This argument also applies to all the finite
matrix solution sets below; it requires no
choice of a countable cofinal sequence.

An actual form over $R$ has a record of some
finite bounds by
Lemma~\ref{C:profinite-projective-data}.
Every element of the split group is a difference
of two actual forms. Pad their idempotent matrices
by zeros and increase the coefficient degree
if necessary to give common bounds $B,D$.
Their reductions satisfy (C.3.PF5).

To prove injectivity, suppose
$[P,q]-[Q,h]$ reduces to zero in every $G_\alpha$.
Over each finite ring, equality of these two
split classes gives an actual equivariant isometry.
Here is the precise cancellation input:
endomorphism rings of finite projective rational
representations over a finite ring are finite
Artinian rings. Their radicals are nilpotent,
idempotents split, and the category has
Krull--Schmidt decomposition.
The hermitian cancellation theorem in
\cite[Theorems 3.2--3.3 and Section 3.4(1),
pp.\ 272--273]{C-QSS1979}, also used in the
finite-base classification, applies.
Thus stable isometry obtained from equality in
the group completion is actual isometry.

Fix two idempotent records of $P,Q$ over $R$.
For each $\alpha$, take all pairs of matrices
$(g_\alpha,u_\alpha)$ solving (C.3.PF3) for the
reduced records. This is a nonempty finite set.
The transition homomorphisms send solutions to
solutions, so the compactness fact gives a
compatible family. Its entries define matrices
$g,u$ over $R$.
All scalar, coaction and Gram identities pass
to the inverse limit: each uses a finite free
coefficient module and finitely many matrix
entries. Hence (C.3.PF3) holds over $R$, giving
an actual equivariant isometry $P\simeq Q$.
This proves injectivity.
In particular we have not assumed cancellation
over the possibly non-Noetherian ring $R$;
it is the consequence of finite cancellation
and these compatible matrix solutions.

Conversely suppose $(\xi_\alpha)$ is compatible
and satisfies (C.3.PF5) with fixed $B,D$.
At each level let $S_\alpha$ be the finite set
of pairs of records of these bounds whose
class difference is $\xi_\alpha$.
It is nonempty by assumption.
Reduction maps $S_\beta$ to $S_\alpha$ whenever
$\beta\geq\alpha$: the record equations are
polynomial, and the required class equality
is preserved by scalar extension.
The compactness fact gives a compatible pair
of records at every level.

Matrices over $R$ are exactly the compatible
matrices over all $R_\alpha$.
The coefficient arrays are also compatible
arrays over $R$, because $C_D$ is finite free
over $\mathbf Z[1/2]$.
Their limits satisfy every equation in
(C.3.PF2). They therefore define actual
finite projective rational symmetric forms
over $R$, with the same ambient matrix bound
$B$ and coefficient bound $D$.
Their difference maps to $(\xi_\alpha)$,
proving the exact image assertion.

The finite decidable test in (C.3.PF5) is
exhaustive enumeration of the records of
Lemma~\ref{C:profinite-projective-data},
followed by the already proved finite-base
class computation and equality test.
The transition maps are computed by reducing
generator records and classifying their images.
This remains an explicit finite procedure for
nonlocal finite rings and projective modules;
it does not replace an ambient idempotent
bound by a free-rank assumption.

The operations are coordinatewise because
base change commutes with them.
Their image stays within (C.3.PF5).
Indeed tensoring two projective presentations
uses $e\otimes f$ in ambient size $BB'$,
and adds coefficient degrees.
A product of differences of forms uses
two direct sums of such tensors, so ambient
size $2BB'$ suffices.
For an exterior power, $\bigwedge^r e$
is an idempotent on
$\bigwedge^r R^B\simeq R^{\binom Br}$,
with image $\bigwedge^r P$.
Its coaction and Gram matrix are determinant
minors, with coefficient degree at most $rD$.
No factorial is divided out.
For a virtual input, expansion of
$\lambda_t(P)\lambda_t(Q)^{-1}$ to a prescribed
order is a finite integral combination of
these tensor and exterior constructions.
Separating positive and negative terms gives
a uniform finite ambient bound depending
only on the input bounds and the specified
order. Thus each requested finite coordinate
is calculated by finite procedures, with
exactly the stated limitation on infinite
data equality.

Finally, the ring (C.3.PF6) is the inverse
limit of the finite rings
$\mathbf F_p\ltimes\mathbf F_p^N$ with
square-zero second summand, under the
surjective coordinate truncations.
Its maximal ideal is the square-zero second
summand, and its residue field is $\mathbf F_p$.
That ideal is not finitely generated:
its multiplication by $R$ factors through
$\mathbf F_p$, while the coordinate vectors
$\epsilon_j$ form an infinite linearly
independent set. Hence this local ring
is not Noetherian. It nonetheless satisfies
every hypothesis of the proved theorem,
which supplies the claimed example beyond
the preceding Noetherian local family.
\end{proof}
